% ====================================================================
% §13 LCO order–disorder 와 MIT 2상역 — 같은 정규용액 틀
% ====================================================================
\section{LCO order--disorder 와 MIT 2상역 --- 같은 정규용액 틀 (N3$'$)}\label{sec:lco-hys}
흑연의 히스테리시스 분기(\S\ref{sec:hys})는 정규용액 이중웰 하나에서 나왔다. 이 절은 LCO 세 전이(T1 MIT$\cdot$
T2$\cdot$T3 order--disorder)를 같은 틀에 넣어 gap 과 분기 중심을 대입형으로 닫고, 정렬 엔트로피와 상호작용
에너지를 혼동하지 않는 경계$\cdot$MIT 의 config/전자 슬롯 분리$\cdot$도핑 보정을 세운다.

\subsection{LCO 세 전이를 같은 정규용액 자유에너지에}\label{sec:lco-hys-gxi}
\textbf{(a) 출발.} Part 0(\S\ref{sec:sm-mf})의 격자기체 화학퍼텐셜 식~\eqref{eq:mu} 와 자유에너지 식~\eqref{eq:gxi}(\S\ref{sec:hys} 가 받아 쓰는)의 유일한
가정 ``동등한 자리에 리튬이 차고 빈다''는 LCO $\mathrm{Li}_x\mathrm{CoO_2}$ 의 팔면체 리튬 자리에도
문자 그대로 성립하므로(자리당 점유 0 또는 1), \S\ref{sec:lco-center} 의 LCO 전이 집합에 번호를 달아
\begin{equation}
\mathcal J_\mathrm{LCO}=\{\mathrm{T1},\mathrm{T2},\mathrm{T3}\}\ \text{(식~\eqref{eq:lco-n0sub})},\qquad
j=1\,(\mathrm{T1{:}\ MIT}),\ \ j=2\,(\mathrm{T2{:}\ OD}_a),\ \ j=3\,(\mathrm{T3{:}\ OD}_b),
\label{eq:lco-J}
\end{equation}
각 전이 $j$ 에 진행률 $\xi_j$ 를 달아 흑연과 같은 정규용액 몫을 쓴다(OD $=$ order--disorder):
\begin{equation}
g_j^\mathrm{cat}(\xi_j)=g_{0,j}^\mathrm{cat}
+RT\{\xi_j\ln\xi_j+(1-\xi_j)\ln(1-\xi_j)\}+\Omega_j^\mathrm{cat}\,\xi_j(1-\xi_j).
\label{eq:lco-gxi}
\end{equation}
새 물리는 넣지 않는다 --- LCO 로 바뀌는 것은 전이별 입력값 $U_j^\mathrm{cat},\Omega_j^\mathrm{cat},
\gamma_j,h_{\eta,j}$ 뿐이다(상첨자 cat 은 같은 슬롯의 LCO 값 표기다).

\textbf{★two-phase calibration.}
흑연에서 spinodal 문턱 $\Omega_j>2RT$($2RT\approx4958$ J/mol@298.15 K, 식~\eqref{eq:spinodal})를 피팅 후
유지할 것으로 지목되는 것은 네 staging 전이 중 $2\mathrm L\!\to\!2$(LiC$_{12}$)$\cdot$$2\!\to\!1$(LiC$_6$)
\emph{두 개} 뿐이고 --- 초기값 $\Omega_j$ 로는 네 전이 모두 문턱 위이나 그 두-상 지목의 실제 근거는 실측
plateau$\cdot$staging 상평형이다(\S\ref{sec:broadening-class}) ---, dilute$\to$stage4 영역($\Omega$ 슬롯이
없는 연속 등온선 --- 전이 아님)과 $4\mathrm L\!\to\!3\mathrm L\cdot3\mathrm L\!\to\!2\mathrm L$ 두 전이(표~\ref{tab:staging}
의 $4\!\to\!3\cdot3\!\to\!2\mathrm L$ 행 --- L 접미는 \S\ref{sec:broadening-class} 의 문헌 명명)는 $\Omega_j<2RT$ 의
solid-solution 쪽으로 피팅될 것이 기대되는(\S\ref{sec:width}$\cdot$\S\ref{sec:broadening}) 반면, LCO 는 pure-LCO
상도표 물리에서 T1(MIT)$\cdot$T2$\cdot$T3(order--disorder) \emph{세 전이 전부}가 같은 문턱
$\Omega_j^\mathrm{cat}>2RT$ 를 넘는 상분리 후보다\cite{reimers1992,vanderven1998,motohashi2009}(표~\ref{tab:lco-staging}). 각 전이의 실제 gap 유무는 최종적으로
\emph{피팅된} $\Omega_j^\mathrm{cat}$ 이 $2RT$ 를 넘는지가 정한다(도핑의 정량형은 \S\ref{sec:lco-hys-dope}·식~\eqref{eq:lco-dope}). 다만 그 후보 지위는 아직 서술
층위다 --- 표~\ref{tab:lco-staging} 는 LCO $\Omega_j^\mathrm{cat}$ 의 수치 열을 싣지 않으며(흑연
표~\ref{tab:staging} 와 달리 초기값 미배정), 미배정 시 $\Omega=0$ 으로 두어 히스 분기가 비활성(0)인
것으로 취급한다. 따라서 본 절의 spinodal$\cdot$gap 식은 $\Omega_j^\mathrm{cat}$
이 round-trip 피팅으로 배정될 때를 위한 구조식이고, 이하의 두-상 서술은 \emph{$\Omega_j^\mathrm{cat}>2RT$
로 피팅되는 전이에 한해} 발효된다.

\subsection{gap 의 닫힌 꼴과 분기 중심 --- 흑연 대입형}\label{sec:lco-hys-gap}
\textbf{(b) 연산 --- 곡률과 spinodal 문턱.}
식~\eqref{eq:lco-gxi} 를 두 번 미분하면 흑연 식~\eqref{eq:gpp} 에 전이별 $\Omega_j^\mathrm{cat}$ 만 든다:
\begin{equation}
\frac{\partial^2 g_j^\mathrm{cat}}{\partial\xi_j^2}=\frac{RT}{\xi_j(1-\xi_j)}-2\Omega_j^\mathrm{cat}.
\label{eq:lco-gpp}
\end{equation}
따라서 spinodal 은
\begin{equation}
\xi_{s,j}^{\pm}=\tfrac12\big(1\pm u_j^\mathrm{cat}\big),\qquad
u_j^\mathrm{cat}=\sqrt{1-\frac{2RT}{\Omega_j^\mathrm{cat}}}\qquad(\Omega_j^\mathrm{cat}>2RT),
\label{eq:lco-spinodal}
\end{equation}
이고 $\Omega_j^\mathrm{cat}<2RT$ 이면 $u_j^\mathrm{cat}$ 이 허수, 등호에선 $u=0$ 이라 그 전이의 spinodal gap 은 없다(gap $0$ 으로 연속 ---
흑연 식~\eqref{eq:dUhys} 와 같은 경우 구분).

\textbf{(c) 중간식 --- LCO 전위 곡선에 spinodal 대입.}
LCO 전이 $j$ 의 평형 전위 곡선은 식~\eqref{eq:Veq} 에 $U_j^\mathrm{cat},\Omega_j^\mathrm{cat}$ 를 넣어
\begin{equation}
V_{\eq,j}^\mathrm{cat}(\xi)=U_j^\mathrm{cat}
+\frac{RT}{F}\ln\frac{\xi}{1-\xi}+\frac{\Omega_j^\mathrm{cat}}{F}(1-2\xi),
\label{eq:lco-Veq}
\end{equation}
이고 두 spinodal 끝점에서 $\dfrac{\xi}{1-\xi}\Big|_{\xi_{s,j}^\pm}=\dfrac{1\pm u_j^\mathrm{cat}}{1\mp
u_j^\mathrm{cat}}$, $(1-2\xi)\big|_{\xi_{s,j}^\pm}=\mp u_j^\mathrm{cat}$ 이므로(흑연 식~\eqref{eq:hyssub}
과 같은 대입) 극대$-$극소 차는(흑연 식~\eqref{eq:hysdiff} 의 전개 재사용)
\[
\Delta U_{j}^{\hys,\mathrm{cat}}=V_{\eq,j}^\mathrm{cat}(\xi_{s,j}^-)-V_{\eq,j}^\mathrm{cat}(\xi_{s,j}^+)
=\frac{2}{F}\big[\Omega_j^\mathrm{cat}u_j^\mathrm{cat}-2RT\,\mathrm{artanh}\,u_j^\mathrm{cat}\big].
\]
상수 중심 $U_j^\mathrm{cat}$ 는 차에서 상쇄되고, \emph{대칭점 검산}으로 두 끝점의 \emph{평균}은 로그 몫과
$(1-2\xi)$ 몫이 각각 부호 반대로 상쇄되어 $\tfrac12[V_{\eq,j}^\mathrm{cat}(\xi_{s,j}^-)
+V_{\eq,j}^\mathrm{cat}(\xi_{s,j}^+)]=U_j^\mathrm{cat}$(흑연 식~\eqref{eq:hyssym} 과 동일 대수) ---
분기가 $U_j^\mathrm{cat}$ 대칭이라는 성질이 LCO 에도 host-무관하게 성립한다.

\textbf{(d) 박스 --- LCO 전이별 gap 과 분기 중심.}
\begin{equation}
\boxed{\;
\Delta U_j^{\hys,\mathrm{cat}}(T)=
\begin{cases}
\dfrac{2}{F}\big[\Omega_j^\mathrm{cat}u_j^\mathrm{cat}-2RT\,\mathrm{artanh}\,u_j^\mathrm{cat}\big],
& \Omega_j^\mathrm{cat}>2RT,\\[4pt]
0, & \Omega_j^\mathrm{cat}\le2RT,
\end{cases}
\quad u_j^\mathrm{cat}=\sqrt{1-\dfrac{2RT}{\Omega_j^\mathrm{cat}}}\;}
\label{eq:lco-dUhys}
\end{equation}
\begin{equation}
\boxed{\;
U_j^{\,d,\mathrm{cat}}=U_j^\mathrm{cat}
+\tfrac12\sigma_d h_{\eta,j}\gamma_j\,\Delta U_j^{\hys,\mathrm{cat}}(T)\;}
\label{eq:lco-Ubranch}
\end{equation}
흑연 식~\eqref{eq:dUhys}$\cdot$\eqref{eq:Ubranch} 에 $\Omega_j\mapsto\Omega_j^\mathrm{cat}$,
$U_j\mapsto U_j^\mathrm{cat}$, $j\in\mathcal J_\mathrm{LCO}$ 를 실제로 넣은 대입형이다. \textbf{★주의 --- $\gamma_j,h_{\eta,j}$ 는 흑연에서와 마찬가지로 여기서도
\emph{유도되지 않는} 방향별 한 자유도의 현상학적 축소 인자다}(두 spinodal 끝점의 대칭 평균
$U_j^\mathrm{cat}$, 식~\eqref{eq:hyssym} 형까지는 엄밀히 유도되나, 그 너머 실측 분기를 한 자유도로 접는
것은 \S\ref{sec:hys} 와 동일 지위의 모델 가정이다).

\textbf{★분기 부호} --- 식~\eqref{eq:lco-Ubranch} 의 $\sigma_d$ 슬롯은 \S\ref{sec:lco-direction} 의 전극-중립
규약(식~\eqref{eq:lco-sigmaslot} --- 탈리튬화 $=+1$, LCO 하프셀에선 충전)을 받고, 가지 배치(탈리튬화
$+\tfrac12$ 위$\cdot$리튬화 $-\tfrac12$ 아래)와 ``탈리튬화 peak 가 높은 전위 쪽'' 읽기는 \S\ref{sec:lco-direction} (b)
항과 동일하게 유지된다.

\subsection{T2$\cdot$T3 order--disorder --- 유효 인력과 $\Omega$/config 구분}\label{sec:lco-hys-od}
$x\approx0.5$ 부근 monoclinic 초격자(점유 Li 와 빈자리의 교대 정렬)를 이루는 T2($\sim$4.05 V,
hex$\to$monoclinic)$\cdot$T3($\sim$4.17--4.20 V, monoclinic$\to$hex)의 실측$\cdot$이론 계보를 먼저
짚는다 --- 이 한 쌍의 좁은 1차 전이와 그 사이 질서상은 in situ X-선 회절로 실측되었고\cite{reimers1992},
제일원리 계산이 $x{=}\tfrac12$ 부근 Li$\cdot$빈자리 질서상의 안정 창과 전이를 재현했으며\cite{vanderven1998},
통계역학–연속체 스케일 브리징 계산이 같은 order--disorder 전이를 연속 조성축 위에서
재확인한다\cite{ml2024}. 이들은 1차 전이의 2상 공존을 전이
진행률 $\xi_j$ 좌표의 정규용액으로 접을 때 나타나는 \emph{유효 인력} $\Omega_j^\mathrm{cat}>0$(미시
기원은 Li$\cdot$빈자리의 정렬 상호작용이고, 이중웰을 만드는 것은 이 유효 좌표의 볼록성 상실이다)이
만드는 상분리의 LCO 사례로, 식~\eqref{eq:lco-dUhys}$\cdot$\eqref{eq:lco-Ubranch} 에 $j=2,3$ 을 넣으면
\[
\begin{aligned}
&u_j^\mathrm{cat}=\sqrt{1-\tfrac{2RT}{\Omega_j^\mathrm{cat}}},\qquad
\Delta U_j^{\hys,\mathrm{cat}}=\tfrac{2}{F}\big[\Omega_j^\mathrm{cat}u_j^\mathrm{cat}
-2RT\,\mathrm{artanh}\,u_j^\mathrm{cat}\big],\\
&U_j^{\,d,\mathrm{cat}}=U_j^\mathrm{cat}+\tfrac12\sigma_d h_{\eta,j}\gamma_j\Delta U_j^{\hys,\mathrm{cat}}
\qquad(j=2,3),
\end{aligned}
\]
로 T2$\cdot$T3 각각 열린다.

\begin{srcbox}
\textbf{다리 --- 유효 인력 $\Omega_j^\mathrm{cat}$ 가 어느 제일원리 상도표에 대응하는가.} T2$\cdot$T3 를
정규용액~\eqref{eq:lco-gxi} 의 \emph{유효 인력} $\Omega_j^\mathrm{cat}>0$ 로 접을 때, 그 미시 기원인
``Li$\cdot$빈자리의 정렬 상호작용''과 $x{=}\tfrac12$ 부근 질서상의 안정성은 Van der Ven
등\cite{vanderven1998} 의 제일원리 상 안정성 계산이 준다. 대응은 다음과 같다(왼쪽이 본문 기호,
오른쪽이 원 논문 량):
\begin{center}\footnotesize
\begin{tabular}{@{}ll@{}}
\hline
본문(식~\eqref{eq:lco-gxi},~\eqref{eq:lco-spinodal}) & Van der Ven\cite{vanderven1998} 대응량 \\
\hline
유효 인력 $\Omega_j^\mathrm{cat}$ (평균장 쌍) & 유효 클러스터 상호작용(ECI) 집합 $\{V_i,V_{ik},\dots\}$ \\
정규용액 볼록성 상실 $\Omega_j^\mathrm{cat}\!>\!2RT$ & 제일원리 자유에너지의 2상 안정역(상분리) \\
$x{\approx}\tfrac12$ 질서상(T2/T3 monoclinic) & 계산이 재현한 $x{=}\tfrac12$ Li$\cdot$빈자리 질서상 \\
저농도 영역(하프셀 상한 밖 옵션) & 저-$x$ staging 화합물(교대 Li 평면 점유) \\
\hline
\end{tabular}
\end{center}
등치의 핵은 \emph{평균장 축약}이다 --- 원 논문의 다체 클러스터 전개를 전이 $j$ 창의 진행률 $\xi_j$
좌표에서 최근접 쌍 한 항으로 접으면, 그 유효 쌍상호작용이 본문의 $\Omega_j^\mathrm{cat}$ 이다:
\begin{equation}
\underbrace{\Omega_j^\mathrm{cat}\,\xi_j(1-\xi_j)}_{\text{식~\eqref{eq:lco-gxi} 정규용액 항}}
\;\longleftrightarrow\;
\underbrace{\big[\text{ECI 전개의 최근접 쌍 몫}\big]_{\text{평균장}}}_{\text{Van der Ven 클러스터 전개의 축약}},
\qquad
\Omega_j^\mathrm{cat}>2RT\ \Leftrightarrow\ \text{$x{\approx}\tfrac12$ 질서상 안정},
\label{eq:br-vanderven1998-1}
\end{equation}
곧 우리 이중웰 문턱 $\Omega_j^\mathrm{cat}>2RT$ 가 켜지는 조건이, 그들의 상도표가 $x{=}\tfrac12$ 질서상을
안정화하는 조성$\cdot$온도 창에 대응한다(통계역학--연속체 스케일 브리징\cite{ml2024} 이 같은 전이를 연속
조성축에서 재확인). \emph{방법 요지} --- Van der Ven 등은 $\mathrm{Li}_x\mathrm{CoO_2}$ 형성에너지를
제일원리 클러스터 전개(자리 점유 변수의 유효 클러스터 상호작용 합)로 적고 이를 상평형(자유에너지
최소)으로 풀어 $x{=}\tfrac12$ Li$\cdot$빈자리 질서상과 저-$x$ staging 을 재현한 상도표를 얻는다.
\emph{가정 차} --- 원 계산은 여러 ECI 와 \emph{명시적} 질서 바닥상태를 전부 담으나, 본문은 이를 전이별
\emph{단일 평균장 쌍상호작용} $\Omega_j^\mathrm{cat}$ 하나로 축약(정규용액 이중웰)하므로 질서상의 미시
구조가 아니라 상분리의 유효 문턱만 남기고, 각 전이의 실제 gap 유무는 최종적으로 \emph{피팅된}
$\Omega_j^\mathrm{cat}$ 이 $2RT$ 를 넘는지가 정한다(pure-LCO $\Omega_j^\mathrm{cat}$ 는 초기값,
\S\ref{sec:lco-hys-dope}).
\end{srcbox}

\textbf{★$\Omega$ 와 config $\Delta S$ 의 구분(혼동 가드)} --- 정렬의 charge-order
엔트로피 변화($\approx0.47$ J/(mol\,K)@$x{=}\tfrac12$, $\approx1.49$ J/(mol\,K)@$x{=}\tfrac23$
[값 원전 재확인 중 --- 종전 귀속처는 전자$\cdot$자기 상도표 논문\cite{motohashi2009}으로 이 $\Delta S$ 값의
1차 원전이 아닌 것으로 확인되어 \emph{값도 배정도} tier C 다. 실측 질서상 조성은
$x{=}\tfrac12$$\cdot$$\tfrac56$\cite{reynier2004} 이어서 $x{=}\tfrac23$ anchor 는 조성 재배정 검토 대상이고 T2/T3
조성 창($x\approx0.55/0.48$)과의 불일치도 그대로다 --- 최종 귀속은 피팅 위임])는 표~\ref{tab:lco-staging}$\cdot$\S\ref{sec:lco-decomp} 의
$\Delta S_j^\mathrm{config}$ 슬롯($U_j^\mathrm{cat}(T)$ 의 온도 이동)에 들어가는 \emph{엔트로피}[J/(mol\,K)]이지
spinodal 문턱을 정하는 \emph{상호작용 에너지} $\Omega_j^\mathrm{cat}$[J/mol]가 \emph{아니므로},
``$0.47/1.49\to\Omega_j^\mathrm{cat}$'' 로 직접 연결하지 않고, $\Omega_j^\mathrm{cat}$ 는 gap 을 정하는
별도 피팅 파라미터로 둔다.

\subsection{T1 MIT 2상역 --- 구조/전자 슬롯 분리}\label{sec:lco-hys-mit}
T1($\sim$3.90 V, $x\approx0.94$--$0.75$, 절연체$\to$금속)의 2상 공존은 in situ 회절\cite{reimers1992}
과 전자 상도표\cite{motohashi2009} 로 실측된 1차 전이이며, 그 미시 기작(Li 빈자리에 속박된 정공
불순물 준위의 Mott 전이\cite{marianetti2004})과 전자 자유도의 배경은 \S\ref{sec:lco-why} 가 다룬다.
이 구조적 2상 공존도 같은 정규용액 문턱을 받아
\[
u_1^\mathrm{cat}=\sqrt{1-\tfrac{2RT}{\Omega_1^\mathrm{cat}}},\quad
\Delta U_1^{\hys,\mathrm{cat}}=\tfrac{2}{F}\big[\Omega_1^\mathrm{cat}u_1^\mathrm{cat}
-2RT\,\mathrm{artanh}\,u_1^\mathrm{cat}\big],\quad
U_1^{\,d,\mathrm{cat}}=U_1^\mathrm{cat}+\tfrac12\sigma_d h_{\eta,1}\gamma_1\Delta U_1^{\hys,\mathrm{cat}}.
\]
MIT 고유의 전자 자유도는 이 히스 gap 에 새 항으로 섞지 않는다 --- 그 항은 이미
$\Delta S_{\rxn,1}^\mathrm{cat}=\Delta S_1^\mathrm{config}+\Delta S_1^\mathrm{vib}
+\Delta S_{e,1}^\mathrm{mol}(x,T)$(\S\ref{sec:lco-decomp} 의 분해식)를 통해 $U_1^\mathrm{cat}(T)$ 의 온도 이동
(식~\eqref{eq:lco-dUdT})에 들어가, 두 몫이 서로 다른 슬롯에 존재한다:
\begin{equation}
\boxed{\;
\underbrace{\Delta U_1^{\hys,\mathrm{cat}}\ \text{(구조적 2상역)}\ \Leftarrow\ \Omega_1^\mathrm{cat}}_{\text{정규용액 히스 gap (이 절)}}
\quad\Big\|\quad
\underbrace{\partial U_1^\mathrm{cat}/\partial T\ \Leftarrow\ \Delta S_{e,1}^\mathrm{mol}(x,T)}_{\text{config$\cdot$vib baseline 에 더해지는 전자 엔트로피 신규 항 (\S\ref{sec:lco-electronic})}}\;}
\label{eq:lco-mit}
\end{equation}
이 분리가 config 2상역과 전자 엔트로피를 이중계산하지 않는 경계다.

\subsection{도핑 보정(우리 시료; 갭 G3)}\label{sec:lco-hys-dope}
Al$^{3+}$/Mg$^{2+}$ 비-redox 치환은 Co redox$\cdot$전자항을 보존하되 order--disorder$\cdot$MIT 상전이를
억제한다 --- 정규용액 틀에서 이는 pure-LCO 상도표 물리\cite{reimers1992,vanderven1998} 의 $\Omega_j^\mathrm{cat,pure}$ 를 도핑 피팅값
$\Omega_j^\mathrm{cat,dop}$ 로 낮추는 입력 변화다. 식~\eqref{eq:lco-dUhys} 의 문턱 극한은
\begin{equation}
\Omega_j^\mathrm{cat,dop}\to2RT^+\ \Longrightarrow\
u_j^\mathrm{cat,dop}\to0,\quad
\Delta U_j^{\hys,\mathrm{cat,dop}}\to\frac{8RT}{3F}\big(u_j^\mathrm{cat,dop}\big)^3\to0,
\label{eq:lco-dope}
\end{equation}
(흑연 식~\eqref{eq:dUhys} 아래 Taylor 극한 재사용, $\mathrm{artanh}\,u=u+u^3/3+\cdots$,
$T_{c,j}=\Omega_j^\mathrm{cat}/2R$)이고, $\Omega_j^\mathrm{cat,dop}\le2RT$ 로 피팅되는 전이는 같은 식 안에서
gap 이 0 으로 닫힌다 --- 이것이 상전이 억제에 따른 히스 축소와 peak smear 의 수식 표현이며, 풀린 몫은
\S\ref{sec:broadening} 의 broadening 폭이 더 크게 담는다. \textbf{★슬롯 분리} --- 중심 전위의 미세 shift 는
같은 전이 파라미터 집합의 $U_j^\mathrm{cat}$ 피팅값 이동으로 \emph{따로} 들어가며, $\Omega_j^\mathrm{cat}$ 하나가
gap$\cdot$smear 와 중심 이동을 동시에 만든다고 쓰지 않는다(흑연 네 전이가 표~\ref{tab:staging} 의 \emph{수치}
초기값에서 출발했던 것과 달리 LCO $\Omega_j^\mathrm{cat}$ 는 수치 초기값 미배정이므로(\S\ref{sec:lco-hys-gxi}),
pure-LCO 상도표 물리는 피팅 배정 시의 \emph{개념적 출발점}이고 폭$\cdot$shift 는 우리 데이터로 피팅).

% 자산: [B2-060] [B2-061] [B2-062] [B2-063] [B2-064] [B2-065] [B2-066] [B2-067] [B2-068] [B2-069] [B2-070] [B2-071] [B2-072] [B2-073] [B2-074]
